% Transcribed from the completed seven-page PDF; proofs and wording preserved.
% Typography correction only: match the reference's 1.12 baseline stretch.
% Build from repository root with XeLaTeX; no shell escape required.
\documentclass[11pt,a4paper]{article}
\usepackage[margin=23mm,headheight=14pt,headsep=8.5mm,footskip=10mm]{geometry}
\usepackage{amsmath,amssymb}
\usepackage[no-math]{fontspec}
\usepackage{xeCJK}
\setmainfont{Latin Modern Roman}
\setCJKmainfont{FandolSong-Regular.otf}[Path=assets/fonts/fandol/,BoldFont=FandolSong-Bold.otf,ItalicFont=FandolSong-Regular.otf,Script=CustomDefault]
\xeCJKsetup{PunctStyle=quanjiao,CheckSingle=true}
\linespread{1.12}
\setlength{\parindent}{2em}\setlength{\parskip}{3pt}
\usepackage{titlesec}
\titleformat{\section}{\centering\normalfont\Large\bfseries}{}{0pt}{}
\titlespacing*{\section}{0pt}{1.8\baselineskip}{1.0\baselineskip}
\usepackage{fancyhdr}\pagestyle{fancy}\fancyhf{}
\fancyhead[L]{\small 复旦大学高等代数习题集·续作}
\fancyfoot[C]{\thepage}\renewcommand{\headrulewidth}{0.4pt}
\usepackage[hidelinks,unicode]{hyperref}
\hypersetup{pdftitle={2022A12—A15、2023S01—S03 七题解答}}
\DeclareMathOperator{\Ker}{Ker}
\DeclareMathOperator{\tr}{tr}
\newcommand{\question}[1]{\fancyhead[R]{\small #1}\section*{问题 #1}}
\newcommand{\sourcepage}[1]{\par\noindent{\small 题目来源：谢启鸿《复旦大学高等代数习题集（2023 版）》，第 #1 页。}}
\newcommand{\proofend}{\hfill$\square$\par}
\widowpenalty=10000\clubpenalty=10000\raggedbottom
\begin{document}
\question{2022A12}
\noindent 设 $V$ 是有理数域上的 $n$ 维线性空间。证明：存在 $V$ 上的 $n$ 个线性变换 $\varphi_1,\ldots,\varphi_n$，使得对任意非零向量 $\alpha\in V$，向量组 $\varphi_1(\alpha),\ldots,\varphi_n(\alpha)$ 均线性无关。

\noindent 证明：设 $n\geq1$。任取 $V$ 的一组基 $e_1,\ldots,e_n$，定义线性变换 $T$：
\[
 T(e_j)=e_{j+1}\quad(1\leq j<n),\qquad T(e_n)=2e_1.
\]
每个基向量经过 $n$ 次作用恰好循环一周，并被乘以 $2$，所以
\[
 T^n=2I_V.
\]
取 $\varphi_j=T^{j-1}$（$1\leq j\leq n$）。下面证明它们满足要求。

先说明 $f(x)=x^n-2$ 在 $\mathbb Q[x]$ 中不可约。$n=1$ 时成立；$n\geq2$ 时，若它在 $\mathbb Q[x]$ 中可约，则由 Gauss 引理可写成两个正次数首一整系数多项式的乘积 $f=uv$。模 $2$ 约化后有 $\overline u\,\overline v=x^n$，故 $\overline u=x^{\deg u}$、$\overline v=x^{\deg v}$。于是 $u(0)$、$v(0)$ 均为偶数，从而 $4\mid u(0)v(0)=-2$，矛盾。这就是本题所需的 Eisenstein 判别法。

这里使用的 Gauss 引理可简述如下：本原整系数多项式的乘积仍为本原多项式，因为模任一素数约化后，两个非零多项式的乘积非零。对有理系数因式清除分母并提取各系数的最大公因数，便得到整系数分解；首一性又使两个因式可取为首一。

任取 $\alpha\ne0$。若有
\[
 c_0\alpha+c_1T\alpha+\cdots+c_{n-1}T^{n-1}\alpha=0,\qquad c_j\in\mathbb Q,
\]
令 $g(x)=c_0+c_1x+\cdots+c_{n-1}x^{n-1}$，则 $g(T)\alpha=0$。若系数不全为零，则 $0\leq\deg g<n$。由 $f$ 不可约，$\gcd(f,g)=1$；辗转相除法给出 $u_1,v_1\in\mathbb Q[x]$，使
\[
 u_1(x)f(x)+v_1(x)g(x)=1.
\]
代入 $T$ 并作用于 $\alpha$，利用 $f(T)=0$，得到
\[
 \alpha=u_1(T)f(T)\alpha+v_1(T)g(T)\alpha=0,
\]
矛盾。因此所有 $c_j$ 均为零，所构造的 $n$ 个向量线性无关。\proofend
\noindent 方法小结：用不可约多项式构造循环变换；低次多项式与它互素，Bézout 恒等式便保证任一非零向量都能生成一组循环基。
\sourcepage{125}

\newpage
\question{2022A13}
\noindent 求一个次数最小的实系数多项式 $f(x)$，使得 $f(x)$ 除以 $x^2+1$ 余 $x+1$，且除以 $x^3+x^2+1$ 余 $x^2-1$。

\noindent 解：记
\[
 a(x)=x^2+1,\qquad b(x)=x^3+x^2+1.
\]
先满足第二个余式条件，设
\[
 f(x)=x^2-1+b(x)h(x).
\]
在模 $a(x)$ 的意义下，$x^2\equiv-1$，因而
\[
 x^2-1\equiv-2,\qquad b(x)\equiv-x.
\]
第一个余式条件遂变为
\[
 -2-xh(x)\equiv x+1\pmod{a(x)},
\]
即 $-xh(x)\equiv x+3$。由于 $(-x)x\equiv1\pmod{a(x)}$，两边乘以 $x$ 得
\[
 h(x)\equiv x(x+3)\equiv3x-1\pmod{a(x)}.
\]
取 $h(x)=3x-1$，便得到
\begin{align*}
 f_0(x)&=x^2-1+(3x-1)(x^3+x^2+1)\\
 &=3x^4+2x^3+3x-2.
\end{align*}
其两个余式也可由下面的等式直接核验：
\begin{align*}
 f_0(x)&=(x^2+1)(3x^2+2x-3)+(x+1),\\
 f_0(x)&=(x^3+x^2+1)(3x-1)+(x^2-1).
\end{align*}
最后证明次数最小。注意到
\[
 xb(x)+(-x^2-x+1)a(x)=1,
\]
所以 $a,b$ 互素。任意另一解 $f$ 与 $f_0$ 具有相同余式，故 $a\mid(f-f_0)$、$b\mid(f-f_0)$，进而 $ab\mid(f-f_0)$。因此所有解恰为
\[
 f(x)=f_0(x)+(x^2+1)(x^3+x^2+1)k(x),\qquad k(x)\in\mathbb R[x].
\]
若 $k\ne0$，右侧第二项的次数至少为 $5$，不能被四次多项式 $f_0$ 的最高次项抵消，所以 $\deg f\geq5$。故唯一的最低次解为
\[
 \boxed{f(x)=3x^4+2x^3+3x-2},\qquad\deg f=4.
\]
\noindent 方法小结：先把一个余式条件写入表达式，再在另一个模多项式下求逆。要证明次数最小，还需说明所有解只相差两个模多项式乘积的倍数。
\sourcepage{125}

\newpage
\question{2022A14}
\noindent 设 $f(x)$ 是有理数域上的不可约多项式，$a$ 与 $1/a$ 都是 $f(x)$ 的复根。证明：若 $b$ 是 $f(x)$ 的任一复根，则 $1/b$ 也是 $f(x)$ 的复根。

\noindent 证明：将 $f$ 除以其首项系数，不改变根与不可约性，故可设
\[
 f(x)=x^n+c_{n-1}x^{n-1}+\cdots+c_1x+c_0,\qquad c_j\in\mathbb Q.
\]
因 $a\ne0$，有 $c_0\ne0$。事实上，若 $c_0=0$，则 $x\mid f$；由 $f$ 不可约且首一，只能有 $f=x$，这与 $f(a)=0$、$a\ne0$ 矛盾。特别地，$f$ 的每个根都非零。

构造倒数多项式
\[
 f^*(x)=x^nf(1/x)=c_0x^n+c_1x^{n-1}+\cdots+c_{n-1}x+1.
\]
由于 $c_0\ne0$，它是 $n$ 次有理系数多项式。题设给出
\[
 f^*(a)=a^nf(1/a)=0,
\]
故 $f$ 与 $f^*$ 有公共根 $a$。

若 $\gcd(f,f^*)=1$，辗转相除法便给出 $u,v\in\mathbb Q[x]$，使
\[
 u(x)f(x)+v(x)f^*(x)=1.
\]
代入 $x=a$ 会得到 $0=1$。因此它们的最大公因式不是常数；又因 $f$ 不可约，必有 $f\mid f^*$。两者次数相同，再比较首项系数，得到
\[
 f^*(x)=c_0f(x),\qquad\text{即}\quad x^nf(1/x)=c_0f(x).
\]
现在任取 $f$ 的根 $b$。由 $b\ne0$，上式给出
\[
 b^nf(1/b)=c_0f(b)=0,
\]
从而 $f(1/b)=0$，结论成立。\proofend
\noindent 补充观察：比较 $f^*=c_0f$ 的常数项，还可得
\[
 1=c_0^2,\qquad c_0=\pm1.
\]
所以题设不仅保证根集在取倒数后保持不变，也限制了首一多项式的常数项。

\noindent 方法小结：把根的变换 $x\mapsto1/x$ 转化为多项式变换 $f\mapsto f^*$。不可约性将“有一个公共根”提升为“两个多项式相差常数倍”，从而控制全部根。
\sourcepage{125}

\newpage
\question{2022A15}
\noindent 证明：实系数多项式 $x^3+px^2+qx+r$ 的三个根（按重数计）都是实数的充要条件是
\[
 4p^3r-p^2q^2-18pqr+4q^3+27r^2\leq0.
\]
\noindent 证明：令 $y=x+p/3$，消去二次项，得
\[
 x^3+px^2+qx+r=H(y):=y^3+Py+Q,
\]
其中
\[
 P=q-\frac{p^2}{3},\qquad Q=r-\frac{pq}{3}+\frac{2p^3}{27}.
\]
实数平移不改变根是否为实数。下面证明 $H$ 的三个根均为实数当且仅当
\[
 4P^3+27Q^2\leq0.
\]
若 $P>0$，则 $H'(y)=3y^2+P>0$，故 $H$ 只有一个实根，且该根为单根；此时 $4P^3+27Q^2>0$。若 $P=0$，则 $H(y)=y^3+Q$；当 $Q\ne0$ 时唯一实根为单根，当 $Q=0$ 时 $0$ 为三重根。因此 $P=0$ 时，三个根均为实数恰好等价于 $Q=0$，也等价于上述不等式。

以下设 $P<0$，记
\[
 t=\sqrt{-P/3}>0,\qquad P=-3t^2.
\]
由 $H'(y)=3(y^2-t^2)$，函数 $H$ 在 $(-\infty,-t)$、$(t,+\infty)$ 上严格递增，在 $(-t,t)$ 上严格递减，且
\[
 H(-t)=Q+2t^3,\qquad H(t)=Q-2t^3.
\]
因 $H(y)\to\pm\infty$（$y\to\pm\infty$），$H$ 有三个实根（按重数计）当且仅当极大值不小于零、极小值不大于零，即
\[
 Q+2t^3\geq0,\qquad Q-2t^3\leq0.
\]
两个不等式都严格时，三个单调区间各有一个实根；取等号时，由
\[
 y^3-3t^2y+2t^3=(y-t)^2(y+2t),\qquad
 y^3-3t^2y-2t^3=(y+t)^2(y-2t)
\]
可知仍有三个实根。若其中一个符号条件不成立，则只剩一个实根。

因此所需条件等价于
\[
 |Q|\leq2t^3\quad\Longleftrightarrow\quad27Q^2\leq108t^6=-4P^3.
\]
最后展开，有
\begin{align*}
 4P^3+27Q^2
 &=4\left(q-\frac{p^2}{3}\right)^3+27\left(r-\frac{pq}{3}+\frac{2p^3}{27}\right)^2\\
 &=4p^3r-p^2q^2-18pqr+4q^3+27r^2.
\end{align*}
于是原结论成立。\proofend
\noindent 方法小结：先平移消去二次项，再比较两个极值与零的大小。严格小于零对应三个互异实根，等于零对应有重根；这样无需直接引用三次方程判别式。
\sourcepage{125}

\newpage
\question{2023S01}
\noindent 设 $V$ 是数域 $K$ 上的 $n$ 维线性空间，$\varphi$ 是 $V$ 上的线性变换。设 $f,g\in K[x]$ 为首一多项式，$d=\gcd(f,g)$、$m=\operatorname{lcm}(f,g)$。证明：
\[
 \Ker d(\varphi)=\Ker f(\varphi)\cap\Ker g(\varphi),\qquad
 \Ker m(\varphi)=\Ker f(\varphi)+\Ker g(\varphi).
\]
\noindent 证明：写成 $f=df_1$、$g=dg_1$，则
\[
 \gcd(f_1,g_1)=1,\qquad m=df_1g_1.
\]
由辗转相除法，存在 $u,v\in K[x]$，满足
\[
 uf_1+vg_1=1,\qquad uf+vg=d.
\]
代入同一线性变换所得的各多项式彼此交换，且乘法关系保持不变。

\noindent （1）最大公因式对应核的交。因 $f(\varphi)=f_1(\varphi)d(\varphi)$、$g(\varphi)=g_1(\varphi)d(\varphi)$，有
\[
 \Ker d(\varphi)\subseteq\Ker f(\varphi)\cap\Ker g(\varphi).
\]
反之，若 $\alpha\in\Ker f(\varphi)\cap\Ker g(\varphi)$，则
\[
 d(\varphi)\alpha=u(\varphi)f(\varphi)\alpha+v(\varphi)g(\varphi)\alpha=0.
\]
这给出反向包含，从而第一个等式成立。

\noindent （2）最小公倍式对应核的和。因为 $f\mid m$、$g\mid m$，所以
\[
 \Ker f(\varphi)+\Ker g(\varphi)\subseteq\Ker m(\varphi).
\]
任取 $\alpha\in\Ker m(\varphi)$。利用 $uf_1+vg_1=1$，将其分解为
\[
 \alpha=\beta+\gamma,\qquad\beta=v(\varphi)g_1(\varphi)\alpha,\qquad
 \gamma=u(\varphi)f_1(\varphi)\alpha.
\]
由 $fg_1=gf_1=m$，分别有
\begin{align*}
 f(\varphi)\beta&=v(\varphi)(fg_1)(\varphi)\alpha=v(\varphi)m(\varphi)\alpha=0,\\
 g(\varphi)\gamma&=u(\varphi)(gf_1)(\varphi)\alpha=u(\varphi)m(\varphi)\alpha=0.
\end{align*}
故 $\beta\in\Ker f(\varphi)$、$\gamma\in\Ker g(\varphi)$，从而 $\alpha\in\Ker f(\varphi)+\Ker g(\varphi)$。第二个等式成立。\proofend
\noindent 方法小结：整除关系给出一个方向的包含，Bézout 恒等式负责反向包含及向量的显式分解。这里一般只是子空间的和；若 $d=1$，交空间为零，才得到直和。
\sourcepage{135}

\newpage
\question{2023S02}
\noindent 设 $V=M_n(\mathbb C)$，将它看作实线性空间。定义
\[
 \varphi(X)=J\overline XJ^{\mathsf T},
\]
其中 $J$ 是基础循环矩阵，$\overline X$ 表示逐项共轭。求 $\varphi$ 的表示矩阵的全体复特征值。

\noindent 解：取实基
\[
 \mathcal B=(E_{11},E_{12},\ldots,E_{nn};\ iE_{11},iE_{12},\ldots,iE_{nn}).
\]
这里 $\dim_{\mathbb R}V=2n^2$。约定下标按模 $n$ 计算，并取 $Je_j=e_{j-1}$，则
\[
 \varphi(E_{ij})=E_{i-1,j-1},\qquad\varphi(iE_{ij})=-iE_{i-1,j-1}.
\]
因此 $\varphi$ 的实表示矩阵为
\[
 M=\begin{pmatrix}S&0\\0&-S\end{pmatrix},
\]
其中 $S$ 表示实矩阵空间上的置换 $E_{ij}\mapsto E_{i-1,j-1}$。

对每个 $d=0,1,\ldots,n-1$，令
\[
 W_d=\operatorname{span}_{\mathbb R}\{E_{i,i+d}:1\leq i\leq n\}.
\]
矩阵单位按差值 $j-i$ 恰好分成这 $n$ 类，故 $M_n(\mathbb R)=\bigoplus_{d=0}^{n-1}W_d$。每个 $W_d$ 都在 $S$ 下不变，且其 $n$ 个基向量构成一个长度为 $n$ 的循环。一个循环的特征多项式为 $\lambda^n-1$：若循环基满足 $Ce_k=e_{k-1}$，则对每个 $\zeta^n=1$，向量 $\sum_{k=0}^{n-1}\zeta^ke_k$ 是特征值 $\zeta$ 的特征向量。它们的特征值互异，故已给出全部 $n$ 个特征值。

于是
\[
 \chi_S(\lambda)=(\lambda^n-1)^n,\qquad
 \chi_{-S}(\lambda)=\bigl(\lambda^n-(-1)^n\bigr)^n,
\]
从而
\[
 \chi_M(\lambda)=(\lambda^n-1)^n\bigl(\lambda^n-(-1)^n\bigr)^n.
\]
据此分两种情形列出答案：

若 $n$ 为奇数，则 $\chi_M(\lambda)=(\lambda^{2n}-1)^n$，全体复特征值为
\[
 \boxed{\lambda_k=e^{k\pi i/n}\quad(0\leq k<2n),\qquad\text{每个的代数重数为}\,n.}
\]
若 $n$ 为偶数，则 $\chi_M(\lambda)=(\lambda^n-1)^{2n}$，全体复特征值为
\[
 \boxed{\lambda_k=e^{2k\pi i/n}\quad(0\leq k<n),\qquad\text{每个的代数重数为}\,2n.}
\]
两种情形的重数总和都为 $2n^2$，与实表示矩阵的阶数一致。由上述循环分解，表示矩阵在 $\mathbb C$ 上可对角化，故相应几何重数也等于代数重数。

\noindent 方法小结：先将复矩阵拆成实部与虚部，共轭使虚部块多出负号；再按矩阵单位的循环轨道计算谱。这里求的是实表示矩阵的复特征值，不能把 $\varphi$ 当作复线性变换。
\sourcepage{135}

\newpage
\question{2023S03}
\noindent 设二阶实方阵 $A,B$ 满足 $AB=BA$。证明：对任意实数 $x,y,z$，
\[
 4xz\det(xA^2+yAB+zB^2)\geq(4xz-y^2)\bigl(x\det A-z\det B\bigr)^2.
\]
\noindent 证明：先说明 $A,B$ 可同时在复数域上化为上三角形。取 $A$ 的一个复特征值 $a_1$，则非零空间 $E=\Ker(A-a_1I)$ 在 $B$ 下不变，因为
\[
 (A-a_1I)Bv=B(A-a_1I)v=0\qquad(v\in E).
\]
$B|_E$ 在复数域上有特征向量，故存在 $A,B$ 的公共特征向量。将它扩充为 $\mathbb C^2$ 的一组基，便有某个可逆复矩阵 $P$ 使
\[
 P^{-1}AP=\begin{pmatrix}a_1&*\\0&a_2\end{pmatrix},\qquad
 P^{-1}BP=\begin{pmatrix}b_1&*\\0&b_2\end{pmatrix}.
\]
记
\[
 \alpha=\det A=a_1a_2,\qquad\beta=\det B=b_1b_2,
\]
以及
\[
 c=a_1b_2+a_2b_1=\tr A\tr B-\tr(AB).
\]
虽然 $a_j,b_j$ 可能不是实数，但 $A,B$ 为实矩阵，故 $\alpha,\beta,c\in\mathbb R$。

令 $D=\det(xA^2+yAB+zB^2)$。由上三角形的对角元可得
\begin{align*}
 D&=(xa_1^2+ya_1b_1+zb_1^2)(xa_2^2+ya_2b_2+zb_2^2)\\
 &=x^2\alpha^2+xy\alpha c+xz(c^2-2\alpha\beta)+y^2\alpha\beta+yz\beta c+z^2\beta^2\\
 &=(x\alpha-z\beta)^2+yc(x\alpha+z\beta)+xzc^2+y^2\alpha\beta.
\end{align*}
因此，原不等式两边之差为
\begin{align*}
 &4xzD-(4xz-y^2)(x\alpha-z\beta)^2\\
 &\quad=4xz\bigl[yc(x\alpha+z\beta)+xzc^2+y^2\alpha\beta\bigr]+y^2(x\alpha-z\beta)^2\\
 &\quad=4x^2z^2c^2+4xyzc(x\alpha+z\beta)+y^2(x\alpha+z\beta)^2\\
 &\quad=\bigl[2xzc+y(x\alpha+z\beta)\bigr]^2\geq0.
\end{align*}
最后一步中的括号为实数，所以原不等式对所有实数 $x,y,z$ 都成立。整个论证没有除以 $x,z$ 或 $4xz-y^2$，也涵盖它们为零的情形。\proofend
\noindent 等号成立当且仅当
\[
 2xz\bigl(\tr A\tr B-\tr(AB)\bigr)+y\bigl(x\det A+z\det B\bigr)=0.
\]
\noindent 方法小结：交换性保证同时上三角化，从而把行列式化为标量乘积；再用迹与行列式把复特征值组合还原为实数，最终得到一个完全平方。
\sourcepage{135}
\end{document}
